\documentclass[../../../include/open-logic-section]{subfiles}

\begin{document}

\olfileid{his}{set}{cantorplane}
\olsection{Cantor论线段与平面}

Schr\"oder-Bernstein（施罗德--伯恩斯坦）定理证明的一些背景，与我们在
\olref[his][set][pathology]{sec} 节中讨论过的历史有关联。
回想一下，1877年，康托尔证明了正方形上的点与其一条边上的点恰好一样多。
在此，我们将给出他第一次尝试的证明。

令 $\unitline$ 为单位线段，即点集 $[0,1]$。
令 $\unitsquare$ 为单位正方形，即点集 $\unitline
\times \unitline$。
用这些记号来说，康托尔 证明了 $\cardeq{\unitline}{\unitsquare}$。
他给 戴德金写了一张便条，实质上包含了如下论证。

\begin{thm}\ollabel{thm:cantorplane}
$\cardeq{\unitline}{\unitsquare}$
\end{thm}

\begin{proof}[证明：第一部分]
任取 $a, b \in \unitline$。将它们写成二进制表示，于是我们得到由 $0$ 和 $1$ 组成的无限序列 $a_1$, $a_2$, \dots, 以及 $b_1$, $b_2$, \dots，使得：
\begin{align*}
a &= 0.a_1a_2a_3a_4\dots\\
b &= 0.b_1b_2b_3b_4\dots
\intertext{现在考虑由下式给出的函数 $f \colon \unitsquare \to \unitline$：}
f(a, b) & = 0.a_1b_1a_2b_2a_3b_3a_4b_4\dots
\end{align*}
现在 $f$ 是一个单射，因为如果 $f(a, b) = f(c,d)$，那么对所有 $n \in \Nat$ 都有 $a_n =
c_n$ 且 $b_n = d_n$，因此 $a = c$ 且 $b = d$。
\end{proof}

遗憾的是，正如 戴德金 向 康托尔 所指出的，这并未解决最初的问题。
考虑 $0.\dot{1}\dot{0} =
0.1010101010\ldots$。我们需要 $f(a,b) = 0.\dot{1}\dot{0}$，其中：
\begin{align*}
a&= 0.\dot{1}\dot{1} = 0.111111\ldots\\
b&= 0
\end{align*}
但是 $a = 0.\dot{1}\dot{1} = 1$。所以，当我们说“将 $a$ 和 $b$ 写成二进制表示”时，我们必须选择使用\emph{哪种}表示法；并且，既然 $f$ 应当是一个\emph{函数}，我们只能使用两种可能表示法中的\emph{一种}。
但是，比如，如果我们使用简单表示法，并将 $a$ 写成“$1.000\ldots$”，那么就不存在任何有序对 $\tuple{a, b}$ 使得 $f(a, b) = 0.\dot{1}\dot{0}$。

总而言之：戴德金 指出，由于存在某些循环小数展开的可能，康托尔 的函数 $f$ 是一个单射但\emph{不是}满射。所以 康托尔 只证明了 $\cardle{\unitsquare}{\unitline}$，而\emph{没有}证明 $\cardeq{\unitsquare}{\unitline}$。

康托尔 几乎立刻回信给 戴德金，实质上建议可以按如下方式完成证明：

\begin{proof}[证明：补全]
这样，我们已经证明了 $\cardle{\unitsquare}{\unitline}$。
但显然存在一个从 $\unitline$ 到 $\unitsquare$ 的单射：只需将线段平放在正方形的一条边上即可。
因此 $\cardle{\unitline}{\unitsquare}$ 且 $\cardle{\unitsquare}{\unitline}$。
根据 Schr\"{o}der--Bernstein 定理（\olref[sfr][siz][sb]{thm:schroder-bernstein}），有 $\cardeq{\unitline}{\unitsquare}$。
\end{proof}

然而，康托尔 当时无法用这些术语写下最后一行，因为 Schr\"{o}der--Bernstein 定理尚未得到证明。
事实上，尽管 康托尔 后来将此表述为一个一般性猜想，但直到1897年它才得到令人满意的证明。（因此，在1877年晚些时候，康托尔 给出了 \olref{thm:cantorplane} 的另一个证明，该证明并未借助 Schr\"{o}der--Bernstein 定理。）

\end{document}